% GENERATED FILE. Edit canonical lessons, then run npm run generate:books.
% canonical-source-hash: fnv1a64:f0e3c17393a75a43
% canonical-lessons: JA-C296-sakan, JA-C296-omairi, JA-C296-maneku, JA-C296-shuukan, JA-C296-uketsugu, JA-R296-review-292-296

\chapter{Customs Still Alive}
\label{ch:ja-customs-still-alive}
\hlchaptermodality{296}

\begin{chapteropening}
\textbf{By the end of this chapter:} \emph{I can say which old customs carry on today.}

Complete the last lesson of chapter 296: I can say which old customs carry on today.
\end{chapteropening}

\section[\texorpdfstring{\ja{さかん} (sakan)}{さかん (sakan)}]{\ja{さかん} (sakan) — thriving; popular}

\label{lesson:JA-C296-sakan}



\begin{quote}
Before the new one: say the Japanese for to ring, then the Japanese for to ring out.
\end{quote}

\subsection*{You'll want to know: \ja{さかん}}
\textbf{\ja{さかん}} — \emph{sakan} — \textquotedblleft{}thriving; popular\textquotedblright{}.

\begin{quote}
\textbf{\ja{今}\ja{も} \ja{その} \ja{まつりは} \ja{さかんです。}} — \emph{ima mo sono matsuri wa sakan desu} — \textquotedblleft{}The festival is still going strong today.\textquotedblright{}
\end{quote}

A na-adjective: \ja{な} before a noun.

\begin{cousinweb}
\textbf{\ja{さかん}} (\emph{sakan}) is a native Japanese word (wago), inherited rather than borrowed from Chinese; its kanji was attached later as a kun'yomi (JMdict).
\end{cousinweb}

\begin{grammarlens}[title={What you've built}]
One more word for telling what happened.
\end{grammarlens}

\subsection*{Your turn}
\begin{itemize}
\raggedright
  \item \emph{Say it:} \emph{sakan}
  \item \emph{Say it:} \emph{sakan}, once more
  \item \emph{Say it:} \emph{ima mo sono matsuri wa sakan desu} — The festival is still going strong today.
  \item \emph{Recall:} say the Japanese for to ring, then the Japanese for to ring out, then say \emph{sakan} again
\end{itemize}

\begin{tcolorbox}[breakable,colback=teal!4,colframe=teal!35!black,title={Before you move on}]
What does \ja{さかん} mean? (\textquotedblleft{}Thriving; popular\textquotedblright{}.) Say it once more.
\end{tcolorbox}
\section[\texorpdfstring{\ja{おまいり} (omairi)}{おまいり (omairi)}]{\ja{おまいり} (omairi) — a visit to a shrine; worship at a temple}

\label{lesson:JA-C296-omairi}



\begin{quote}
Before the new one: say the Japanese for to ring out, then the Japanese for thriving.
\end{quote}

\subsection*{You'll want to know: \ja{おまいり}}
\textbf{\ja{おまいり}} — \emph{omairi} — \textquotedblleft{}a visit to a shrine; worship at a temple\textquotedblright{}.

\begin{quote}
\textbf{\ja{まいつき} \ja{じんじゃに} \ja{おまいりに} \ja{いきます。}} — \emph{maitsuki jinja ni omairi ni ikimasu} — \textquotedblleft{}Every month I go to the shrine to pray.\textquotedblright{}
\end{quote}

\begin{cousinweb}
\textbf{\ja{おまいり}} (\emph{omairi}) is a native Japanese word (wago), inherited rather than borrowed from Chinese, built from \emph{o-}, \textquotedblleft{}the polite prefix\textquotedblright{}, and \emph{mairi}, \textquotedblleft{}going (humbly)\textquotedblright{} (from \emph{mairu}) (JMdict, KANJIDIC2).
\end{cousinweb}

\begin{grammarlens}[title={What you've built}]
One more word for telling what happened.
\end{grammarlens}

\subsection*{Your turn}
\begin{itemize}
\raggedright
  \item \emph{Say it:} \emph{omairi}
  \item \emph{Say it:} \emph{omairi}, once more
  \item \emph{Say it:} \emph{maitsuki jinja ni omairi ni ikimasu} — Every month I go to the shrine to pray.
  \item \emph{Recall:} say the Japanese for to ring out, then the Japanese for thriving, then say \emph{omairi} again
\end{itemize}

\begin{tcolorbox}[breakable,colback=teal!4,colframe=teal!35!black,title={Before you move on}]
What does \ja{おまいり} mean? (\textquotedblleft{}A visit to a shrine; worship at a temple\textquotedblright{}.) Say it once more.
\end{tcolorbox}
\section[\texorpdfstring{\ja{まねく} (maneku)}{まねく (maneku)}]{\ja{まねく} (maneku) — to beckon; to invite}

\label{lesson:JA-C296-maneku}



\begin{quote}
Before the new one: say the Japanese for thriving, then the Japanese for a visit to a shrine.
\end{quote}

\subsection*{You'll want to know: \ja{まねく}}
\textbf{\ja{まねく}} — \emph{maneku} — \textquotedblleft{}to beckon; to invite\textquotedblright{}.

\begin{quote}
\textbf{\ja{まつりに} \ja{ともだちを} \ja{まねきました。}} — \emph{matsuri ni tomodachi o manekimashita} — \textquotedblleft{}I asked friends to the festival.\textquotedblright{}
\end{quote}

Polite \textbf{\ja{まねきます}}, \emph{manekimasu}; a godan verb. Past (polite): \textbf{\ja{ちちが} \ja{ともだちを} \ja{まねきました。}}, \emph{chichi ga tomodachi o manekimashita}.

\begin{cousinweb}
\textbf{\ja{まねく}} (\emph{maneku}) is a native Japanese word (wago), inherited rather than borrowed from Chinese; its kanji was attached later as a kun'yomi (JMdict).
\end{cousinweb}

\begin{grammarlens}[title={What you've built}]
One more word for telling what happened.
\end{grammarlens}

\subsection*{Your turn}
\begin{itemize}
\raggedright
  \item \emph{Say it:} \emph{maneku}
  \item \emph{Say it:} \emph{maneku}, once more
  \item \emph{Say it:} \emph{matsuri ni tomodachi o manekimashita} — I asked friends to the festival.
  \item \emph{Recall:} say the Japanese for thriving, then the Japanese for a visit to a shrine, then say \emph{maneku} again
\end{itemize}

\begin{tcolorbox}[breakable,colback=teal!4,colframe=teal!35!black,title={Before you move on}]
What does \ja{まねく} mean? (\textquotedblleft{}To beckon; to invite\textquotedblright{}.) Say it once more.
\end{tcolorbox}
\section[\texorpdfstring{\ja{しゅうかん} (shūkan)}{しゅうかん (shūkan)}]{\ja{しゅうかん} (shūkan) — a custom; a habit}

\label{lesson:JA-C296-shuukan}



\begin{quote}
Before the new one: say the Japanese for a visit to a shrine, then the Japanese for to beckon.
\end{quote}

\subsection*{You'll want to know: \ja{しゅうかん}}
\textbf{\ja{しゅうかん}} — \emph{shūkan} — \textquotedblleft{}a custom; a habit\textquotedblright{}.

\begin{quote}
\textbf{\ja{日本}\ja{には} \ja{おもしろい} \ja{しゅうかんが} \ja{あります。}} — \emph{nihon ni wa omoshiroi shūkan ga arimasu} — \textquotedblleft{}Japan has some interesting customs.\textquotedblright{}
\end{quote}

\begin{cousinweb}
\textbf{\ja{しゅうかん}} (\emph{shūkan}) is Sino-Japanese (kango): two Chinese characters read with their on'yomi, readings borrowed from Chinese. They are \emph{shū}, \textquotedblleft{}learn\textquotedblright{} (Middle Chinese \emph{zip}) and \emph{kan}, \textquotedblleft{}accustomed\textquotedblright{} (\emph{kwaenH}) (KANJIDIC2; Middle Chinese in Baxter's transcription).
\end{cousinweb}

\begin{grammarlens}[title={What you've built}]
One more word for telling what happened.
\end{grammarlens}

\subsection*{Your turn}
\begin{itemize}
\raggedright
  \item \emph{Say it:} \emph{shūkan}
  \item \emph{Say it:} \emph{shūkan}, once more
  \item \emph{Say it:} \emph{nihon ni wa omoshiroi shūkan ga arimasu} — Japan has some interesting customs.
  \item \emph{Recall:} say the Japanese for a visit to a shrine, then the Japanese for to beckon, then say \emph{shūkan} again
\end{itemize}

\begin{tcolorbox}[breakable,colback=teal!4,colframe=teal!35!black,title={Before you move on}]
What does \ja{しゅうかん} mean? (\textquotedblleft{}A custom; a habit\textquotedblright{}.) Say it once more.
\end{tcolorbox}
\section[\texorpdfstring{\ja{うけつぐ} (uketsugu)}{うけつぐ (uketsugu)}]{\ja{うけつぐ} (uketsugu) — to inherit; to carry on}

\label{lesson:JA-C296-uketsugu}



\begin{quote}
Before the new one: say the Japanese for to beckon, then the Japanese for a custom.
\end{quote}

\subsection*{You'll want to know: \ja{うけつぐ}}
\textbf{\ja{うけつぐ}} — \emph{uketsugu} — \textquotedblleft{}to inherit; to carry on\textquotedblright{}.

\begin{quote}
\textbf{\ja{あには} \ja{ちちの} \ja{しごとを} \ja{うけつぎました。}} — \emph{ani wa chichi no shigoto o uketsugimashita} — \textquotedblleft{}My brother took over our father's work.\textquotedblright{}
\end{quote}

Polite \textbf{\ja{うけつぎます}}, \emph{uketsugimasu}; a godan verb. Past (polite): \textbf{\ja{むすこが} \ja{みせを} \ja{うけつぎました。}}, \emph{musuko ga mise o uketsugimashita}.

\begin{cousinweb}
\textbf{\ja{うけつぐ}} (\emph{uketsugu}) is a native Japanese word (wago), inherited rather than borrowed from Chinese, built from \emph{uke}, \textquotedblleft{}receiving\textquotedblright{} (from \emph{ukeru}), and \emph{tsugu}, \textquotedblleft{}to succeed to\textquotedblright{} (JMdict, KANJIDIC2).
\end{cousinweb}

\begin{grammarlens}[title={What you've built}]
One more word for telling what happened.
\end{grammarlens}

\subsection*{Your turn}
\begin{itemize}
\raggedright
  \item \emph{Say it:} \emph{uketsugu}
  \item \emph{Say it:} \emph{uketsugu}, once more
  \item \emph{Say it:} \emph{ani wa chichi no shigoto o uketsugimashita} — My brother took over our father's work.
  \item \emph{Recall:} say the Japanese for to beckon, then the Japanese for a custom, then say \emph{uketsugu} again
\end{itemize}

\begin{tcolorbox}[breakable,colback=teal!4,colframe=teal!35!black,title={Before you move on}]
What does \ja{うけつぐ} mean? (\textquotedblleft{}To inherit; to carry on\textquotedblright{}.) Say it once more.
\end{tcolorbox}
\section[(dialogue)]{Looking back: from A Long Life to Customs Still Alive}

\label{lesson:JA-R296-review-292-296}



\begin{quote}
Say the Japanese for to grow old and an old person.
\end{quote}

\subsection*{Your turn}
Say these aloud:

\begin{itemize}
\raggedright
  \item to grow old, an old person, a life, to die, to support
  \item fireworks, to set off (fireworks), a food stall, the night sky, to bustle
  \item a portable shrine, to carry on one's shoulders, a procession, to liven up, a Bon festival dance
  \item the New Year shrine visit, a shrine fortune slip, a charm, to ring, to ring out
  \item thriving, a visit to a shrine, to beckon, a custom, to inherit
\end{itemize}

\begin{tcolorbox}[breakable,colback=teal!4,colframe=teal!35!black,title={Before you move on}]
To grow old? (\textbf{\ja{おいる}}, \emph{oiru}.) An old person? (\textbf{\ja{としより}}, \emph{toshiyori}.) A life? (\textbf{\ja{じんせい}}, \emph{jinsei}.) To die? (\textbf{\ja{しぬ}}, \emph{shinu}.) To support? (\textbf{\ja{ささえる}}, \emph{sasaeru}.) Fireworks? (\textbf{\ja{はなび}}, \emph{hanabi}.) To set off (fireworks)? (\textbf{\ja{うちあげる}}, \emph{uchiageru}.) A food stall? (\textbf{\ja{やたい}}, \emph{yatai}.) The night sky? (\textbf{\ja{よぞら}}, \emph{yozora}.) To bustle? (\textbf{\ja{にぎわう}}, \emph{nigiwau}.) A portable shrine? (\textbf{\ja{みこし}}, \emph{mikoshi}.) To carry on one's shoulders? (\textbf{\ja{かつぐ}}, \emph{katsugu}.) A procession? (\textbf{\ja{ぎょうれつ}}, \emph{gyōretsu}.) To liven up? (\textbf{\ja{もりあがる}}, \emph{moriagaru}.) A Bon festival dance? (\textbf{\ja{ぼんおどり}}, \emph{bon'odori}.) The New Year shrine visit? (\textbf{\ja{はつもうで}}, \emph{hatsumōde}.) A shrine fortune slip? (\textbf{\ja{おみくじ}}, \emph{omikuji}.) A charm? (\textbf{\ja{おまもり}}, \emph{omamori}.) To ring? (\textbf{\ja{ならす}}, \emph{narasu}.) To ring out? (\textbf{\ja{ひびく}}, \emph{hibiku}.) Thriving? (\textbf{\ja{さかん}}, \emph{sakan}.) A visit to a shrine? (\textbf{\ja{おまいり}}, \emph{omairi}.) To beckon? (\textbf{\ja{まねく}}, \emph{maneku}.) A custom? (\textbf{\ja{しゅうかん}}, \emph{shūkan}.) To inherit? (\textbf{\ja{うけつぐ}}, \emph{uketsugu}.)
\end{tcolorbox}
